% GENERATED FILE. Edit canonical lessons, then run npm run generate:books.
% canonical-source-hash: fnv1a64:80be69926c9fc1c8
% canonical-lessons: KA-C409-prabandha, KA-C409-sanksepisu, KA-C409-grantha, KA-C409-mandisu, KA-C409-tippani

\chapter{Writing an Essay}
\label{ch:ka-writing-an-essay}
\hlchaptermodality{409}

\begin{chapteropening}
\textbf{By the end of this chapter:} \emph{I can say an essay, to summarise, a treatise, to present and a note.}

Complete the last lesson of chapter 409: I can say an essay, to summarise, a treatise, to present and a note.
\end{chapteropening}

\section[\texorpdfstring{\kn{ಪ್ರಬಂಧ} (prabandha)}{ಪ್ರಬಂಧ (prabandha)}]{\kn{ಪ್ರಬಂಧ} (prabandha) — an essay}

\label{lesson:KA-C409-prabandha}



\begin{quote}
Before the new one: say the Kannada for a syllabus, then the Kannada for revision.
\end{quote}

\subsection*{You'll want to know: \kn{ಪ್ರಬಂಧ}}
\textbf{\kn{ಪ್ರಬಂಧ}} — \emph{prabandha} — \textquotedblleft{}an essay\textquotedblright{}.

\begin{quote}
\kn{ನಾನು} \kn{ನಮ್ಮ} \kn{ಊರಿನ} \kn{ಬಗ್ಗೆ} \kn{ಒಂದು} \kn{ಪ್ರಬಂಧ} \kn{ಬರೆದೆ}.
\end{quote}

\begin{quote}
\emph{nānu namma ūrina bagge ondu prabandha barede.}
\end{quote}

\begin{quote}
I wrote an essay about our town.
\end{quote}

Also a thesis or a paper.

\begin{cousinweb}
\textbf{\kn{ಪ್ರಬಂಧ}} (\emph{prabandha}) is borrowed from Sanskrit \emph{prabandha}, \textquotedblleft{}a connection, a continuous composition\textquotedblright{} (CDIAL 8695).
\end{cousinweb}

\begin{grammarlens}[title={What you've built}]
One more word for giving reasons and opinions.
\end{grammarlens}

\subsection*{Your turn}
\begin{itemize}
\raggedright
  \item \emph{Say it:} \emph{prabandha}
  \item \emph{Say it:} \emph{prabandha}, once more
  \item \emph{Say it:} \emph{nānu namma ūrina bagge ondu prabandha barede.} — I wrote an essay about our town.
  \item \emph{Recall:} say the Kannada for a syllabus, then the Kannada for revision, then say \emph{prabandha} again
\end{itemize}

\begin{tcolorbox}[breakable,colback=teal!4,colframe=teal!35!black,title={Before you move on}]
What does \kn{ಪ್ರಬಂಧ} mean? (\textquotedblleft{}An essay\textquotedblright{}.) Say it once more.
\end{tcolorbox}
\section[\texorpdfstring{\kn{ಸಂಕ್ಷೇಪಿಸು} (saṅkṣēpisu)}{ಸಂಕ್ಷೇಪಿಸು (saṅkṣēpisu)}]{\kn{ಸಂಕ್ಷೇಪಿಸು} (saṅkṣēpisu) — to summarise; to abridge}

\label{lesson:KA-C409-sanksepisu}



\begin{quote}
Before the new one: say the Kannada for revision, then the Kannada for an essay.
\end{quote}

\subsection*{You'll want to know: \kn{ಸಂಕ್ಷೇಪಿಸು}}
\textbf{\kn{ಸಂಕ್ಷೇಪಿಸು}} — \emph{saṅkṣēpisu} — \textquotedblleft{}to summarise; to abridge\textquotedblright{}.

\begin{quote}
\kn{ನಾನು} \kn{ಹತ್ತು} \kn{ಪುಟದ} \kn{ಲೇಖನವನ್ನು} \kn{ಒಂದು} \kn{ಪುಟಕ್ಕೆ} \kn{ಸಂಕ್ಷೇಪಿಸಿದೆ}.
\end{quote}

\begin{quote}
\emph{nānu hattu puṭada lēkhanavannu ondu puṭakke saṅkṣēpiside.}
\end{quote}

\begin{quote}
I cut a ten-page article down to one page.
\end{quote}

Past (I …): \kn{ಸಂಕ್ಷೇಪಿಸಿದೆ}, \emph{saṅkṣēpiside}.

\kn{ಸಂಕ್ಷಿಪ್ತ}, brief; \kn{ಸಂಕ್ಷಿಪ್ತವಾಗಿ}, in brief.

\begin{cousinweb}
\textbf{\kn{ಸಂಕ್ಷೇಪಿಸು}} (\emph{saṅkṣēpisu}) is built on Sanskrit \emph{saṃkṣepa}, \textquotedblleft{}compression, abridgment\textquotedblright{} (Monier-Williams), with the Kannada verb ending -\emph{isu}.
\end{cousinweb}

\begin{grammarlens}[title={What you've built}]
One more word for giving reasons and opinions.
\end{grammarlens}

\subsection*{Your turn}
\begin{itemize}
\raggedright
  \item \emph{Say it:} \emph{saṅkṣēpisu}
  \item \emph{Say it:} \emph{saṅkṣēpisu}, once more
  \item \emph{Say it:} \emph{nānu hattu puṭada lēkhanavannu ondu puṭakke saṅkṣēpiside.} — I cut a ten-page article down to one page.
  \item \emph{Recall:} say the Kannada for revision, then the Kannada for an essay, then say \emph{saṅkṣēpisu} again
\end{itemize}

\begin{tcolorbox}[breakable,colback=teal!4,colframe=teal!35!black,title={Before you move on}]
What does \kn{ಸಂಕ್ಷೇಪಿಸು} mean? (\textquotedblleft{}To summarise; to abridge\textquotedblright{}.) Say it once more.
\end{tcolorbox}
\section[\texorpdfstring{\kn{ಗ್ರಂಥ} (grantha)}{ಗ್ರಂಥ (grantha)}]{\kn{ಗ್ರಂಥ} (grantha) — a treatise; a book}

\label{lesson:KA-C409-grantha}



\begin{quote}
Before the new one: say the Kannada for an essay, then the Kannada for to summarise.
\end{quote}

\subsection*{You'll want to know: \kn{ಗ್ರಂಥ}}
\textbf{\kn{ಗ್ರಂಥ}} — \emph{grantha} — \textquotedblleft{}a treatise; a book\textquotedblright{}.

\begin{quote}
\kn{ಅಜ್ಜನ} \kn{ಮನೆಯಲ್ಲಿ} \kn{ತುಂಬಾ} \kn{ಹಳೆಯ} \kn{ಗ್ರಂಥಗಳು} \kn{ಇವೆ}.
\end{quote}

\begin{quote}
\emph{ajjana maneyalli tumbā haḷeya granthagaḷu ive.}
\end{quote}

\begin{quote}
Grandfather's house has many old books.
\end{quote}

Plural: \kn{ಗ್ರಂಥಗಳು}, \emph{granthagaḷu}.

A weighty or old book. \kn{ಗ್ರಂಥಾಲಯ}, a library.

\begin{cousinweb}
\textbf{\kn{ಗ್ರಂಥ}} (\emph{grantha}) is borrowed from Sanskrit \emph{grantha}, \textquotedblleft{}a knot\textquotedblright{}, and later \textquotedblleft{}a book\textquotedblright{} (CDIAL 4350), from \emph{granth}, \textquotedblleft{}to tie\textquotedblright{}.
\end{cousinweb}

\begin{grammarlens}[title={What you've built}]
One more word for giving reasons and opinions.
\end{grammarlens}

\subsection*{Your turn}
\begin{itemize}
\raggedright
  \item \emph{Say it:} \emph{grantha}
  \item \emph{Say it:} \emph{grantha}, once more
  \item \emph{Say it:} \emph{ajjana maneyalli tumbā haḷeya granthagaḷu ive.} — Grandfather's house has many old books.
  \item \emph{Recall:} say the Kannada for an essay, then the Kannada for to summarise, then say \emph{grantha} again
\end{itemize}

\begin{tcolorbox}[breakable,colback=teal!4,colframe=teal!35!black,title={Before you move on}]
What does \kn{ಗ್ರಂಥ} mean? (\textquotedblleft{}A treatise; a book\textquotedblright{}.) Say it once more.
\end{tcolorbox}
\section[\texorpdfstring{\kn{ಮಂಡಿಸು} (maṇḍisu)}{ಮಂಡಿಸು (maṇḍisu)}]{\kn{ಮಂಡಿಸು} (maṇḍisu) — to present}

\label{lesson:KA-C409-mandisu}



\begin{quote}
Before the new one: say the Kannada for to summarise, then the Kannada for a treatise.
\end{quote}

\subsection*{You'll want to know: \kn{ಮಂಡಿಸು}}
\textbf{\kn{ಮಂಡಿಸು}} — \emph{maṇḍisu} — \textquotedblleft{}to present\textquotedblright{}.

\begin{quote}
\kn{ನಾನು} \kn{ಸಮ್ಮೇಳನದಲ್ಲಿ} \kn{ಒಂದು} \kn{ಪ್ರಬಂಧ} \kn{ಮಂಡಿಸಿದೆ}.
\end{quote}

\begin{quote}
\emph{nānu sammēḷanadalli ondu prabandha maṇḍiside.}
\end{quote}

\begin{quote}
I presented a paper at the conference.
\end{quote}

Past (I …): \kn{ಮಂಡಿಸಿದೆ}, \emph{maṇḍiside}.

For a paper, a case or a bill. \kn{ಮಂಡನೆ}, a presentation.

\begin{cousinweb}
\textbf{\kn{ಮಂಡಿಸು}} (\emph{maṇḍisu}) is a Kannada verb whose origin is not settled: neither DEDR nor CDIAL lists it.
\end{cousinweb}

\begin{grammarlens}[title={What you've built}]
One more word for giving reasons and opinions.
\end{grammarlens}

\subsection*{Your turn}
\begin{itemize}
\raggedright
  \item \emph{Say it:} \emph{maṇḍisu}
  \item \emph{Say it:} \emph{maṇḍisu}, once more
  \item \emph{Say it:} \emph{nānu sammēḷanadalli ondu prabandha maṇḍiside.} — I presented a paper at the conference.
  \item \emph{Recall:} say the Kannada for to summarise, then the Kannada for a treatise, then say \emph{maṇḍisu} again
\end{itemize}

\begin{tcolorbox}[breakable,colback=teal!4,colframe=teal!35!black,title={Before you move on}]
What does \kn{ಮಂಡಿಸು} mean? (\textquotedblleft{}To present\textquotedblright{}.) Say it once more.
\end{tcolorbox}
\section[\texorpdfstring{\kn{ಟಿಪ್ಪಣಿ} (ṭippaṇi)}{ಟಿಪ್ಪಣಿ (ṭippaṇi)}]{\kn{ಟಿಪ್ಪಣಿ} (ṭippaṇi) — a note; a comment}

\label{lesson:KA-C409-tippani}



\begin{quote}
Before the new one: say the Kannada for a treatise, then the Kannada for to present.
\end{quote}

\subsection*{You'll want to know: \kn{ಟಿಪ್ಪಣಿ}}
\textbf{\kn{ಟಿಪ್ಪಣಿ}} — \emph{ṭippaṇi} — \textquotedblleft{}a note; a comment\textquotedblright{}.

\begin{quote}
\kn{ಉಪನ್ಯಾಸದ} \kn{ಸಮಯದಲ್ಲಿ} \kn{ನಾನು} \kn{ಟಿಪ್ಪಣಿ} \kn{ಮಾಡಿಕೊಂಡೆ}.
\end{quote}

\begin{quote}
\emph{upanyāsada samayadalli nānu ṭippaṇi māḍikoṇḍe.}
\end{quote}

\begin{quote}
I took notes during the lecture.
\end{quote}

\kn{ಟಿಪ್ಪಣಿ} \kn{ಮಾಡಿಕೊ}, to take notes for oneself.

\begin{cousinweb}
\textbf{\kn{ಟಿಪ್ಪಣಿ}} (\emph{ṭippaṇi}) is borrowed from Sanskrit \emph{ṭippaṇī}, \textquotedblleft{}a gloss, a commentary\textquotedblright{} (Monier-Williams).
\end{cousinweb}

\begin{grammarlens}[title={What you've built}]
One more word for giving reasons and opinions.
\end{grammarlens}

\subsection*{Your turn}
\begin{itemize}
\raggedright
  \item \emph{Say it:} \emph{ṭippaṇi}
  \item \emph{Say it:} \emph{ṭippaṇi}, once more
  \item \emph{Say it:} \emph{upanyāsada samayadalli nānu ṭippaṇi māḍikoṇḍe.} — I took notes during the lecture.
  \item \emph{Recall:} say the Kannada for a treatise, then the Kannada for to present, then say \emph{ṭippaṇi} again
\end{itemize}

\begin{tcolorbox}[breakable,colback=teal!4,colframe=teal!35!black,title={Before you move on}]
What does \kn{ಟಿಪ್ಪಣಿ} mean? (\textquotedblleft{}A note; a comment\textquotedblright{}.) Say it once more.
\end{tcolorbox}
